% Source: https://artofproblemsolving.com/community/c3h2600125_insert_something_creative_here
% Stable identity: aops-topic-2600125-q01
\textbf{(a)} There is exactly one way to choose no objects (choose the empty set), so $\dbinom n0=1$. For the formula to give this we need
\[1=\dfrac{n!}{0!\,n!}=\dfrac1{0!},\]
so $0!=1$. The recursion $1!=1\cdot0!$ also forces $0!=1!=1$. So $0!=1$ is not proved from the definition $n!=1\cdot2\cdots n$ (which only makes sense for $n\ge1$); it is the only value that keeps these formulas true, and that is why it is chosen.

\textbf{(b)} Each $2$-element subset $\{i,j\}$ with $i<j$ has a larger element $j\in\{2,3,\dots,n+1\}$. For a given $j$ there are $j-1$ choices of $i$, namely $1,\dots,j-1$. Adding over $j$,
\[\dbinom{n+1}{2}=\sum_{j=2}^{n+1}(j-1)=1+2+\cdots+n.\ \blacksquare\]
(This is why the triangular numbers $1,3,6,10,\dots$ appear in Pascal's triangle as $\dbinom22,\dbinom32,\dbinom42,\dbinom52,\dots$)
